\section{Критерий Сильвестра}

\begin{definition}
Пусть $A$ --- симметричная матрица квадратичной формы $f(x, x)$ порядка $n$. \textbf{Главными минорами} $\Delta_k$ матрицы $A$ называются определители её верхних левых квадратных подматриц порядка $k$ ($k = 1, \dots, n$):
\[
\Delta_1 = a_{11}, \quad 
\Delta_2 = \begin{vmatrix} a_{11} & a_{12} \\ a_{21} & a_{22} \end{vmatrix}, \quad \dots, \quad 
\Delta_n = \det A = \begin{vmatrix} 
a_{11} & a_{12} & \dots & a_{1n} \\ 
a_{21} & a_{22} & \dots & a_{2n} \\ 
\vdots & \vdots & \ddots & \vdots \\ 
a_{n1} & a_{n2} & \dots & a_{nn} 
\end{vmatrix}.
\]
По договоренности полагаем $\Delta_0 = 1$.
\end{definition}

\begin{theorem}[Критерий Сильвестра]
Квадратичная форма $f(x, x)$ с симметричной матрицей $A$ в некотором базисе является:
\begin{enumerate}
    \item \textbf{Положительно определенной} тогда и только тогда, когда все её главные миноры строго положительны:
    \[
    \Delta_1 > 0, \quad \Delta_2 > 0, \quad \dots, \quad \Delta_n > 0.
    \]
    \item \textbf{Отрицательно определенной} тогда и только тогда, когда знаки главных миноров строго чередуются, начиная с отрицательного:
    \[
    \Delta_1 < 0, \quad \Delta_2 > 0, \quad \Delta_3 < 0, \quad \dots, \quad (-1)^n \Delta_n > 0.
    \]
\end{enumerate}
\end{theorem}

\begin{proof}
\begin{enumerate}
    \item \textbf{Необходимость (для положительной определенности):} 
    Предположим, что $f(x, x) > 0$ для всех $x \neq 0$, но существует главный минор $\Delta_k = 0$ для некоторого $k \le n$. Рассмотрим однородную систему линейных уравнений, составленную из первых $k$ уравнений системы $Ax = 0$:
    \[
    \begin{cases}
    a_{11}x_1 + a_{12}x_2 + \dots + a_{1k}x_k = 0 \\
    a_{21}x_1 + a_{22}x_2 + \dots + a_{2k}x_k = 0 \\
    \vdots \\
    a_{k1}x_1 + a_{k2}x_2 + \dots + a_{kk}x_k = 0
    \end{cases}
    \]
    Определитель этой системы равен $\Delta_k = 0$, следовательно, она имеет нетривиальное решение $x_0 = (x_1^0, \dots, x_k^0, 0, \dots, 0)^T \neq 0$. 
    Вычислим значение квадратичной формы на этом векторе:
    \[
    f(x_0, x_0) = \sum_{i=1}^n \sum_{j=1}^n a_{ij} x_i^0 x_j^0 = \sum_{i=1}^k x_i^0 \left( \sum_{j=1}^k a_{ij} x_j^0 \right).
    \]
    Поскольку $x_0$ является решением системы, выражение в скобках равно нулю для каждого $i = 1, \dots, k$. Следовательно, $f(x_0, x_0) = 0$. Это противоречит положительной определенности формы (так как $x_0 \neq 0$, должно быть $f(x_0, x_0) > 0$). 
    Таким образом, все $\Delta_k \neq 0$. Поскольку все главные миноры ненулевые, к форме применим метод Якоби, приводящий её к каноническому виду:
    \[
    f(x, x) = \frac{\Delta_1}{\Delta_0} y_1^2 + \frac{\Delta_2}{\Delta_1} y_2^2 + \dots + \frac{\Delta_n}{\Delta_{n-1}} y_n^2.
    \]
    Чтобы эта форма была положительно определенной, все коэффициенты при квадратах должны быть строго положительны: $\frac{\Delta_k}{\Delta_{k-1}} > 0$ для всех $k$. Поскольку $\Delta_0 = 1 > 0$, методом математической индукции получаем $\Delta_1 > 0, \Delta_2 > 0, \dots, \Delta_n > 0$.

    \item \textbf{Достаточность:} 
    Если все $\Delta_k > 0$, то метод Якоби применим, и форма приводится к каноническому виду с коэффициентами $\frac{\Delta_k}{\Delta_{k-1}} > 0$. Следовательно, для любого $x \neq 0$ (а значит, и $y \neq 0$) сумма квадратов с положительными коэффициентами строго больше нуля. Форма положительно определена.

    \item \textbf{Отрицательная определенность:} 
    Форма $f(x, x)$ отрицательно определена тогда и только тогда, когда форма $-f(x, x)$ положительно определена. Матрица формы $-f$ равна $-A$. Её главные миноры равны $(-1)^k \Delta_k$. Применяя к $-f$ критерий положительной определенности, получаем условие $(-1)^k \Delta_k > 0$ для всех $k$, что эквивалентно знакопеременности миноров $\Delta_k$, начиная с отрицательного.
\end{enumerate}
\hfill$\blacksquare$
\end{proof}

\begin{remark}
Критерий Сильвестра позволяет однозначно классифицировать формы на положительно и отрицательно определенные. Однако он не различает случаи полуопределенности и знакопеременности (для этого требуется анализ всех, а не только главных миноров, или использование индексов инерции).
\end{remark}
